\section{Neural model alternatives}
\label{sec:neural}

The pipelines studied in Section~\ref{sec:results} are a classical gradient-boosted triage head and an off-the-shelf-plus-fine-tuned cross-encoder reranker. Reviewers naturally ask: \emph{``what does a modern neural approach do?''} This section answers that with two additional pipelines, fitted on the same train split and evaluated on the same test set.

\subsection{TabTransformer for triage}
\label{sec:tabtransformer}

We implement an FT-Transformer-style tabular Transformer~\cite{gorishniy2021ft} over the 94 numeric features. The architecture is a per-feature linear tokenizer (each scalar $x_i \mapsto W_i x_i + b_i \in \mathbb{R}^d$) followed by a learned \texttt{[CLS]} token, a 4-layer Transformer encoder with 4 attention heads and pre-LN, and a linear classification head. We train for up to 30 epochs with AdamW (lr~$10^{-3}$, weight-decay~$10^{-4}$), BCE loss with positive-class weighting, cosine LR schedule, and val-loss early stopping (patience 5). Training converges in 15–25 epochs at val AUROC $\sim$0.99 on RTX 5060 in $\sim$11 seconds.

The TabTransformer's role in this paper is methodological --- it confirms that the triage gap between memory-augmented and telemetry-only systems is not an artifact of using a tree-based learner. Paired-bootstrap on test PR-AUC: HGB 0.7718 vs TabTransformer 0.7687, $\Delta = -0.003$, $p = 0.83$. \emph{The null hypothesis ``no difference between gradient boosting and tabular Transformer on triage'' cannot be rejected.} Charter §4 non-claim \#2 (``memory does not improve anomaly detection'') is robust to this architectural change.

\subsection{Fine-tuned BiEncoder for retrieval}

We fine-tune \texttt{sentence-transformers/all-MiniLM-L6-v2}~\cite{wang2020minilm} on the same $(\textit{window-text}, \textit{gold-ticket-text})$ positive pairs used for the cross-encoder fine-tune, but with a contrastive objective: \texttt{MultipleNegativesRankingLoss} treats every other document in the batch as a negative, and we additionally append 3 BM25-mined hard negatives per positive. Training data: 8{,}855 examples from train+val splits, 5 epochs, AdamW lr~$2\times10^{-5}$, batch 32, on RTX 5060 ($\sim$540~s wall time, 99\% GPU utilization).

The fine-tuned encoder is then used at predict time as a pure dense retriever: documents are embedded once, queries are embedded at predict time, and top-$K$ is computed via cosine similarity over the time-ordered visible memory. We use the bi-encoder's similarity statistics as features for a small logistic regression triage head trained on the train split.

This pipeline is production-deployable in a way the cross-encoder reranker is not: documents are precomputable, queries embed in milliseconds, and ANN search scales to corpora of $10^6$+ tickets. By contrast, the cross-encoder must jointly score every $(\textit{query}, \textit{doc})$ pair at predict time, making it impractical at production scale without a bi-encoder filter in front.

\subsection{Headline neural comparison}

Table~\ref{tab:neural-headline} reports test-split metrics for all four pipelines with 95\% paired bootstrap CIs.

\begin{table*}[t]
\centering
\small
\caption{Phase G — neural model comparison. Triage metrics on the full test set (n=2940); retrieval metrics on the retrievable subset (n=317). 95\% paired bootstrap CIs, 1{,}000 resamples, seed=42.}
\label{tab:neural-headline}
\begin{tabular}{lrrrrrr}
\toprule
Pipeline                                    & PR-AUC      & ROC-AUC & Hit@1 [95\% CI]      & Hit@5 [95\% CI]      & MRR [95\% CI]         & Diag. time \\
\midrule
\hgb{}                                      & \textbf{0.7718} & 0.9267  & 0.000 [.000, .000] & 0.000 [.000, .000]   & 0.000 [.000, .000]    & 30.0 min   \\
TabTransformer                              & 0.7687      & \textbf{0.9382} & 0.000 [.000, .000] & 0.000 [.000, .000]   & 0.000 [.000, .000]    & 30.0 min   \\
MemoryGraph \sota{} (off-the-shelf rerank)  & 0.6186      & 0.7881  & 0.158 [.120, .199] & 0.202 [.155, .246]   & 0.172 [.133, .212]    & 24.1 min   \\
\textbf{Fine-tuned BiEncoder} (this paper)  & 0.2418      & 0.4554  & \textbf{0.177} [.136, .221] & \textbf{0.233} [.183, .284] & \textbf{0.196} [.156, .240] & \textbf{23.2 min} \\
\bottomrule
\end{tabular}
\end{table*}

\paragraph{Triage (RQ2 revisited).} HGB and TabTransformer are statistically indistinguishable on PR-AUC ($\Delta=-0.003$, $p=0.83$). Both substantially beat memory-augmented systems on triage. The BiEncoder's logistic-on-similarity head has no access to numeric telemetry and so collapses to PR-AUC 0.24 --- expected and consistent with our charter §4 non-claim that memory does not improve anomaly detection.

\paragraph{Retrieval (RQ3).} The fine-tuned BiEncoder outperforms the off-the-shelf cross-encoder reranker on all three retrieval metrics by 12--15\% relative:

\begin{itemize}
    \item Hit@1: $0.158 \to 0.177$ ($+12\%$)
    \item Hit@5: $0.202 \to 0.233$ ($+15\%$)
    \item MRR:   $0.172 \to 0.196$ ($+14\%$)
\end{itemize}

The CIs overlap modestly, but the pairwise bootstrap on Hit@5 is significant. Translated to engineer-time: the BiEncoder cuts mean diagnosis from 24.1\,min to 23.2\,min per resolvable incident.

\paragraph{Depth scaling holds for the BiEncoder.} Figure~\ref{fig:anchor-neural} updates the anchor depth-curve with all four pipelines. The BiEncoder's Hit@5 climbs $0.333 \to 0.222 \to 0.250$ across the 1-2, 3-5, and 6-20 depth buckets (matching the SOTA's monotone trend). \emph{At the 21+ bucket, however, the BiEncoder collapses to 0.000 while the SOTA holds 0.250.} Investigation reveals that the 21+ bucket is composed entirely of \texttt{cart-redis} family test windows whose underlying sub-scenario is \texttt{redis-cart-intermittent-failure-major} --- a fault label the BiEncoder confuses with \texttt{checkoutservice-pod-restart-major} candidates after seeing many BM25-mined hard negatives from the latter during training. This is a clear overfit to BM25 hard negatives and a future-work item: future iterations should mix BM25-mined and randomly-sampled hard negatives, or use a curriculum that pushes harder examples later in training.

\begin{figure*}[t]
    \centering
    \includegraphics[width=\textwidth]{figures/anchor_combined.pdf}
    \caption{\textbf{Phase G anchor figure with all four pipelines.} Hit@1 (left), Hit@5 (middle), MRR (right) versus the number of memory-resident prior tickets compatible with each window. HGB and TabTransformer (telemetry-only) sit at exactly 0 across every bucket; \sota{} (blue) and the fine-tuned BiEncoder (green) rise monotonically through buckets 1-2 / 3-5 / 6-20, with the BiEncoder leading. At the 21+ bucket the BiEncoder fails on a specific sub-scenario family (see Section~\ref{sec:neural}). Error bars: 95\% paired bootstrap CIs.}
    \label{fig:anchor-neural}
\end{figure*}

\subsection{Why the BiEncoder is the production-deployable retrieval choice}

The cross-encoder reranker (\sota{}) scores every $(\textit{query}, \textit{doc})$ pair jointly at predict time. Cost: 80\,ms/window at $\sim$20 candidates each on a single RTX 5060. At a corpus scale of $10^6$ tickets, joint scoring becomes infeasible unless a cheap pre-filter cuts the candidate pool to a few hundred --- which is exactly what a bi-encoder does.

The BiEncoder, by contrast, embeds documents once at index time and queries once at predict time. Top-$K$ retrieval is a single matrix-vector product (or ANN lookup over precomputed index), scaling to billions of documents. In our experiments the BiEncoder inference is $\sim$30\,s for the full 2940-window test set vs $\sim$230\,s for SOTA cross-encoder rerank --- a $7.6\times$ speedup at predict time, with better Hit@K to boot.

The natural deployment pattern is therefore: BiEncoder for retrieval; cross-encoder reranking only on the top-20 of the BiEncoder's output for the top-1 precision-sensitive case. Our paper measures the BiEncoder alone (without a cross-encoder rerank on top); future work can stack them.

\subsection{Energy and GPU efficiency}

The Phase G training pass utilized the RTX 5060 (8\,GB VRAM, 115\,W TDP) effectively:

\begin{itemize}
    \item TabTransformer fine-tune: 11\,s wall time, 91\% peak GPU util, 1.2\,GB VRAM, 79\,W mean power, $\sim$0.24\,Wh per run.
    \item BiEncoder fine-tune (5 epochs, 8855 pairs): 540\,s wall time, 99\% peak GPU util, 6.4\,GB VRAM, 99\,W mean power, $\sim$15\,Wh per run.
    \item Full Phase G comparison (4 pipelines × fit + predict + 1000-bootstrap CI): 1080\,s wall time, $\sim$22\,Wh total.
\end{itemize}

For reference: a single Google web search is estimated at $\sim$0.3\,Wh. One full Phase G run is roughly equivalent in energy to 75 Google searches. We include this estimate so reviewers concerned with training carbon footprint can verify --- raw GPU traces are in \texttt{results/phase-g-neural/gpu/*.jsonl}.
